\documentclass[11pt,a4paper]{article}

\usepackage[utf8]{inputenc}
\usepackage[T1]{fontenc}
\usepackage[spanish,es-tabla,es-nodecimaldot,es-noquoting]{babel}
\usepackage{lmodern}
\usepackage{microtype}
\usepackage[margin=2.4cm]{geometry}
\usepackage{graphicx}
\usepackage{xcolor}
\usepackage{tabularx,booktabs,array}
\usepackage{enumitem}
\usepackage{float}
\usepackage{caption}
\usepackage{fancyhdr}
\usepackage{titlesec}
\usepackage{listings}
\usepackage[most]{tcolorbox}
\usepackage{tikz}
\usetikzlibrary{positioning,arrows.meta,fit,backgrounds,calc,babel}
\usepackage[hidelinks]{hyperref}
\usepackage{xurl}

% Mismos tokens que el tema oscuro de la app
\definecolor{noche}{HTML}{0F1726}
\definecolor{panel}{HTML}{1B2638}
\definecolor{ok}{HTML}{2E9E6A}
\definecolor{ambar}{HTML}{E0A526}
\definecolor{alerta}{HTML}{D64545}
\definecolor{tinta}{HTML}{18202C}
\definecolor{gris}{HTML}{667085}
\definecolor{linea}{HTML}{D9DEE7}
\definecolor{fondo}{HTML}{F3F5F9}

\graphicspath{{img/}}

\titleformat{\section}{\Large\bfseries\color{noche}}{\thesection.}{0.6em}{}
\titleformat{\subsection}{\large\bfseries\color{tinta}}{\thesubsection.}{0.6em}{}
\titleformat{\subsubsection}{\normalsize\bfseries\color{gris}}{}{0em}{}

\pagestyle{fancy}
\fancyhf{}
\fancyhead[L]{\small\color{gris}Vigía · HMI/HRI para Unitree Go2}
\fancyhead[R]{\small\color{gris}Desarrollo de Aplicaciones I · UADE}
\fancyfoot[C]{\small\thepage}
\renewcommand{\headrulewidth}{0.4pt}

\setlist{itemsep=2pt,topsep=4pt}
\setlength{\parskip}{0.5em}
\setlength{\parindent}{0pt}
\renewcommand{\arraystretch}{1.25}
\newcolumntype{L}{>{\raggedright\arraybackslash}X}

\lstset{basicstyle=\ttfamily\footnotesize,frame=single,rulecolor=\color{linea},
  backgroundcolor=\color{fondo},xleftmargin=4pt,xrightmargin=4pt,columns=fullflexible,
  breaklines=true,
  literate={á}{{\'a}}1 {é}{{\'e}}1 {í}{{\'i}}1 {ó}{{\'o}}1 {ú}{{\'u}}1 {ñ}{{\~n}}1 {ü}{{\"u}}1}

\newtcolorbox{nota}[1][]{enhanced,colback=fondo,colframe=ambar,arc=2mm,
  boxrule=0.6pt,left=6pt,right=6pt,#1}

% Dato que el equipo todavía tiene que completar o verificar
\newcommand{\pendiente}[1]{\textcolor{red!70!black}{[#1]}}

\begin{document}

\begin{titlepage}
\pagecolor{noche}\color{white}
\renewcommand{\pendiente}[1]{\textcolor{ambar}{[#1]}}
\vspace*{1.5cm}
{\large Universidad Argentina de la Empresa\par Facultad de Ingeniería y Ciencias Exactas\par}
\vspace{0.4cm}
{\large Desarrollo de Aplicaciones I · Trabajo Práctico Grupal\par}

\vspace{4cm}
{\fontsize{46}{50}\selectfont\bfseries Vigía\par}
\vspace{0.3cm}
{\LARGE\color{ambar} La ronda nocturna, supervisada desde la central\par}
\vspace{0.8cm}
{\Large Interfaz HMI/HRI para un Unitree Go2 con telemetría simulada\par}

\vfill
\begin{tabular}{@{}ll@{}}
\textbf{Integrantes} & Papa, Valentin · LU 1197189 \\
                     & Velazquez, Joaquin Eduardo · LU 1186723 \\
                     & Jerez, Juan · LU 1152497 \\
                     & Maci, Abril · LU 1152146 \\
                     & \pendiente{confirmar integrantes de este TP} \\
\textbf{Docente}     & \pendiente{completar} \\
\textbf{Fecha}       & \pendiente{fecha de entrega} \\
\end{tabular}
\end{titlepage}
\nopagecolor\color{black}

\tableofcontents
\newpage

\section{Problema y servicio}
\label{sec:problema}

De noche la sede de la facultad queda con un solo guardia de seguridad, de 22:00 a 06:00. Ese guardia tiene dos maneras de hacer la ronda y las dos tienen un costo.

Si sale a caminar el piso, la central queda vacía: nadie mira las cámaras ni atiende el teléfono mientras dura el recorrido. Si se queda en la central, depende de cámaras fijas que no cubren todo. Los pasillos laterales, el depósito y la salida de emergencia son puntos ciegos conocidos. En ningún caso queda un registro confiable de qué se revisó y a qué hora; a la mañana siguiente el jefe de seguridad solo tiene la palabra del guardia y lo que haya grabado alguna cámara.

\subsection{Qué proponemos}

Un Unitree Go2 hace la ronda por una ruta fija de puntos de control en planta baja. En cada punto se detiene unos segundos y lo marca como verificado. Si encuentra una puerta abierta, algo que le bloquea el paso o una persona en una zona restringida, avisa a la central y espera la decisión del guardia.

El guardia no maneja el robot con un joystick. Lo supervisa desde una app Android en la central y le da órdenes concretas: iniciar una ronda, pausarla, mandarlo a la base, reconocer una alerta. La app se llama Vigía.

\subsection{Cómo se resuelve hoy y por qué un robot}

\begin{table}[H]
\centering\small
\begin{tabularx}{\textwidth}{@{}l L L@{}}
\toprule
\textbf{Alternativa} & \textbf{Qué resuelve} & \textbf{Qué no resuelve} \\
\midrule
Guardia que camina & Ve todo de cerca & Deja la central sola; no hay registro de puntos verificados \\
Más cámaras fijas & Cubren más superficie desde la central & Siguen teniendo ángulos muertos y no pueden acercarse a revisar una puerta \\
Segundo guardia nocturno & Uno camina y otro mira & Es el costo más alto y es justamente lo que la institución no quiere sumar \\
Go2 haciendo la ronda & El guardia se queda en la central y el robot camina; cada punto queda registrado con hora & Necesita supervisión, carga de batería y un gateway propio (ver sección~\ref{sec:robot}) \\
\bottomrule
\end{tabularx}
\end{table}

Este uso no es una idea nuestra. Hay empresas de seguridad que ya ofrecen rondas con robots cuadrúpedos en plantas industriales y depósitos, con un operador humano que decide ante cada evento. \pendiente{agregar una fuente de un caso real}

\section{Elección del robot}
\label{sec:robot}

\subsection{Go2 frente a G1}

La ronda pide caminar varias horas por pasillos, pasar umbrales y algún escalón bajo, y mirar alrededor. No pide agarrar ni mover objetos. Para eso el Go2 alcanza y sobra: es más chico, más estable sobre cuatro patas y bastante más barato que el G1. Los brazos del humanoide no se usarían, y un robot de altura humana caminando de noche por un pasillo asusta más que un cuadrúpedo del tamaño de un perro grande.

\subsection{Qué trae el Go2}

Unitree vende el Go2 en tres versiones: AIR, PRO y EDU. Para este servicio haría falta la EDU, que es la única que habilita el SDK para mandarle órdenes desde un sistema propio.

\begin{table}[H]
\centering\small
\begin{tabularx}{\textwidth}{@{}l L@{}}
\toprule
\textbf{Aspecto} & \textbf{Lo que usamos del Go2} \\
\midrule
Sensores & LiDAR 4D en la cabeza y cámara frontal. Con eso detecta obstáculos y arma el mapa del piso. \\
Movilidad & Camina en interiores, sube escalones bajos y se recupera si lo empujan. Para la ronda limitamos la velocidad a 0,8~m/s. \\
Autonomía & Del orden de una a dos horas según la batería y el uso \pendiente{verificar con la hoja técnica}. Una ronda completa dura unos 15 minutos, así que entran varias por carga. \\
Comunicación & WiFi. La app no le habla directo: pasa por un gateway propio en la red del edificio. \\
Interacción & Parlante y luces en algunas versiones, que usamos para avisar a las personas que se cruza. \\
\bottomrule
\end{tabularx}
\end{table}

\subsection{Qué existe y qué habría que desarrollar}

\begin{table}[H]
\centering\small
\begin{tabularx}{\textwidth}{@{}L L@{}}
\toprule
\textbf{Disponible} & \textbf{Requiere equipamiento o desarrollo extra} \\
\midrule
Locomoción, mapeo y evitación de obstáculos & Detectar si una puerta quedó abierta (modelo de visión propio) \\
SDK de la versión EDU para órdenes de alto nivel & Gateway entre la app y el robot, con autenticación y registro \\
Cámara y LiDAR para detectar que hay alguien & Decidir si esa persona está o no autorizada (eso lo decide el guardia) \\
Batería intercambiable & Base de carga automática, que se compra a terceros o se reemplaza por un cambio manual de batería \\
\bottomrule
\end{tabularx}
\end{table}

En el TP no hay hardware. Todo lo que el robot ``hace'' lo hace un simulador dentro de la app, y el contrato de mensajes que proponemos (sección~\ref{sec:contrato}) es nuestro, no el protocolo oficial de Unitree.

\section{Entorno y recorrido}
\label{sec:entorno}

El robot trabaja en la planta baja de la sede, de 22:00 a 06:00. Para el TP usamos un plano simplificado de 40 por 20 metros con ocho puntos de control y una base de carga al lado de la central de seguridad. Las coordenadas del plano son las mismas que usa el simulador de la app.

\begin{figure}[H]
\centering
\begin{tikzpicture}[x=0.38cm,y=0.38cm,
  zona/.style={draw=linea,fill=fondo,thick},
  restr/.style={draw=alerta,fill=alerta!8,thick,dashed},
  cp/.style={circle,draw=noche,fill=white,thick,inner sep=1.4pt,font=\scriptsize\bfseries},
  lbl/.style={font=\scriptsize,text=gris,align=center}]
  % Zonas
  \draw[zona] (0,0) rectangle (6,6);   \node[lbl] at (3,5) {Central};
  \draw[zona] (6,0) rectangle (14,8);  \node[lbl] at (10,7) {Hall};
  \draw[zona] (14,3) rectangle (26,7); \node[lbl] at (20,6.2) {Pasillo A};
  \draw[zona] (14,7) rectangle (26,12);\node[lbl] at (20,9.5) {Aulas};
  \draw[restr] (26,0) rectangle (36,8);\node[lbl] at (31,7) {Laboratorio};
  \draw[zona] (30,8) rectangle (40,16);\node[lbl] at (36.5,15) {Biblioteca};
  \draw[restr] (26,14) rectangle (34,20);\node[lbl] at (29,19) {Sala técnica};
  \draw[zona] (6,12) rectangle (26,16);\node[lbl] at (14,15.2) {Pasillo B};
  \draw[zona] (16,16) rectangle (24,20);\node[lbl] at (20,19.2) {Salida de emergencia};
  % Ruta completa
  \draw[-{Stealth[length=2mm]},ambar,very thick]
    (3,3) -- (8,2) -- (12,5) -- (20,5) -- (30,4) -- (34,12) -- (30,17) -- (20,18) -- (10,14) -- (3,3);
  % Puntos de control
  \node[cp,fill=ok,text=white] at (3,3) {B};
  \node[cp] at (8,2) {1};
  \node[cp] at (12,5) {2};
  \node[cp] at (20,5) {3};
  \node[cp,draw=alerta] at (30,4) {4};
  \node[cp] at (34,12) {5};
  \node[cp,draw=alerta] at (30,17) {6};
  \node[cp] at (20,18) {7};
  \node[cp] at (10,14) {8};
\end{tikzpicture}
\caption{Plano simplificado de la planta baja con la ronda completa. En rojo punteado, las zonas restringidas.}
\label{fig:plano}
\end{figure}

\begin{table}[H]
\centering\small
\begin{tabularx}{\textwidth}{@{}l L l c@{}}
\toprule
\textbf{Id} & \textbf{Punto de control} & \textbf{Tipo} & \textbf{Posición (m)} \\
\midrule
BASE & Base de carga, junto a la central & Punto seguro & (3, 3) \\
CP1 & Hall de entrada y puerta principal & Puerta & (8, 2) \\
CP2 & Molinetes de acceso & Paso & (12, 5) \\
CP3 & Pasillo de aulas A & Paso & (20, 5) \\
CP4 & Laboratorio de informática & Zona restringida & (30, 4) \\
CP5 & Biblioteca & Paso & (34, 12) \\
CP6 & Sala técnica & Zona restringida & (30, 17) \\
CP7 & Salida de emergencia lateral & Puerta & (20, 18) \\
CP8 & Pasillo de aulas B & Paso & (10, 14) \\
\bottomrule
\end{tabularx}
\end{table}

Hay tres rutas cargadas. La \textbf{ronda completa} pasa por los ocho puntos y vuelve a la base. La \textbf{ronda perimetral} revisa solo accesos y puertas (CP1, CP7 y CP2). La de \textbf{zonas críticas} va directo al laboratorio y a la sala técnica, y sirve para cuando el guardia quiere una revisión rápida de lo más sensible.

\section{Actores e interacción}
\label{sec:interaccion}

\subsection{Roles}

\begin{table}[H]
\centering\small
\begin{tabularx}{\textwidth}{@{}l L L@{}}
\toprule
\textbf{Rol} & \textbf{Qué necesita saber} & \textbf{Cuándo interviene} \\
\midrule
Guardia de la central & Dónde está el robot, qué puntos verificó, si hay alertas abiertas y si los datos están al día & Inicia la ronda, decide ante un incidente, reconoce alertas, manda el robot a la base \\
Jefe de seguridad & Qué rondas se hicieron, cuándo y qué pasó en cada una & Al día siguiente, revisando el historial \\
Personas en el edificio & Qué es el robot y qué tienen que hacer cuando lo cruzan & Cuando se lo encuentran en un pasillo o en una zona restringida \\
\bottomrule
\end{tabularx}
\end{table}

\subsection{HMI: lo que ve el guardia}

La pantalla principal es la de supervisión. Arriba está el estado de la conexión con la hora del último dato; en el centro, el plano con la ruta, los puntos ya verificados y la posición del robot; abajo, las órdenes disponibles. Cada orden muestra si se puede usar y, si no, por qué.

Usamos tema oscuro por defecto. El guardia trabaja de noche en una central con poca luz y una pantalla blanca encandila. Los colores tienen un solo significado en toda la app: verde para lo normal, ámbar para advertencias, rojo para errores y gris para datos viejos.

\subsection{HRI: lo que ven las personas que se cruzan con el robot}

El robot no está pensado para hablar con la gente, pero se la va a cruzar: personal de limpieza, un docente que se quedó tarde, alguien de mantenimiento. Para esos casos definimos un comportamiento fijo:

\begin{itemize}
  \item Si detecta a una persona en un lugar de paso, se detiene a dos metros, le da paso y sigue.
  \item Si la persona está en una zona restringida, se detiene a distancia y dice por el parlante: ``Zona restringida. Seguridad fue notificada''. No la sigue ni la encierra.
  \item El robot nunca toma decisiones sobre la persona. Avisa y es el guardia quien la clasifica.
\end{itemize}

\subsection{Cuándo tiene que intervenir un humano}

El robot sigue la ronda solo, pero frena y espera al guardia en cuatro situaciones: algo le bloquea el camino, detecta a una persona en una zona restringida, la batería baja del umbral crítico o pierde la conexión por más de un minuto. En este último caso no puede esperar a nadie, así que vuelve solo a la base.

\section{Telemetría y órdenes}
\label{sec:telemetria}

\subsection{Datos que muestra la app}

\begin{table}[H]
\centering\small
\begin{tabularx}{\textwidth}{@{}l L@{}}
\toprule
\textbf{Dato} & \textbf{Cómo se muestra} \\
\midrule
Conectividad & Conectado, señal débil o sin datos, con la latencia en milisegundos \\
Batería & Porcentaje y minutos estimados. Ámbar debajo del 25~\%, rojo debajo del 15~\% \\
Estado operativo & En base, listo, patrullando, en pausa, incidente, regresando, detenido o falla \\
Misión y progreso & Ruta activa, punto de control actual y porcentaje recorrido \\
Ubicación & Posición sobre el plano y nombre de la zona \\
Movimiento & Velocidad y rumbo \\
Alertas & Cantidad de alertas activas y la más grave \\
Última actualización & Hora del último dato y su antigüedad. Con más de 3~s sin datos, todo pasa a gris con la leyenda ``hace 12~s'' \\
\bottomrule
\end{tabularx}
\end{table}

Las tres variables propias del servicio son:

\begin{enumerate}
  \item \textbf{Puntos de control verificados} en la ronda actual, por ejemplo 3 de 8. Es lo que el jefe de seguridad quiere ver al día siguiente.
  \item \textbf{Personas detectadas en zonas restringidas} durante la ronda.
  \item \textbf{Puertas abiertas detectadas}, que de noche tendrían que estar todas cerradas.
\end{enumerate}

\subsection{Órdenes}

Cada orden pasa por cuatro estados posibles, y la app siempre muestra en cuál está: \textit{pendiente} mientras espera respuesta, \textit{confirmada} si el robot la aceptó, \textit{rechazada} si la rechazó con un motivo, y \textit{sin respuesta} si pasaron 5~segundos sin noticias.

\begin{nota}
Que una orden esté confirmada quiere decir que el robot la aceptó, no que la terminó. Que el robot efectivamente se detuvo o llegó a la base se ve después, en la telemetría. Y la detención de la app es un comando de software simulado: no reemplaza la parada de emergencia física del robot.
\end{nota}

\begin{table}[H]
\centering\small
\begin{tabularx}{\textwidth}{@{}l L l@{}}
\toprule
\textbf{Orden} & \textbf{Se habilita cuando} & \textbf{Rechazo típico} \\
\midrule
Iniciar misión & Hay datos recientes, el robot está en base o listo, la batería está en 30~\% o más y no hay un incidente abierto & Batería baja, robot ocupado \\
Pausar & Está patrullando & Estado inválido \\
Reanudar & Está en pausa, con 15~\% de batería o más y sin incidente abierto & Batería baja \\
Cancelar misión & Patrullando, en pausa, en incidente o detenido. Pide confirmación & Estado inválido \\
Regresar a punto seguro & En cualquier estado salvo en base o ya regresando. Pide confirmación & Camino bloqueado \\
Detención operativa & Siempre. Hay que mantener apretado el botón 1,5~s & Puede quedar sin respuesta \\
Reconocer alerta & Hay una alerta activa sin reconocer & Alerta desconocida \\
Resolver incidente & Hay un incidente abierto. El guardia elige la decisión y después la confirma & Alerta desconocida \\
\bottomrule
\end{tabularx}
\end{table}

Mientras una orden está pendiente no se puede mandar otra del mismo tipo. La excepción es la detención operativa: se puede volver a mandar siempre, aunque haya una pendiente. Sin datos recientes los botones siguen visibles y habilitados, porque el guardia tiene que poder intentar detener el robot, pero la app avisa que la orden puede no llegar. Solo iniciar misión se deshabilita, porque necesita datos al día.

Cancelar se permite también con el robot detenido. Si no, después de una detención la única salida sería mandarlo a la base.

\subsection{Contrato de mensajes propuesto}
\label{sec:contrato}

Este contrato es una propuesta nuestra para un gateway propio entre la app y el robot. No es el protocolo oficial de Unitree. El repositorio simulado lo usa de verdad: serializa cada orden a JSON y parsea la telemetría, las respuestas y las alertas que le devuelve el simulador.

Los ejemplos que siguen son mensajes reales del simulador, tomados de una ronda de prueba. Los \texttt{timestamp} viajan en ISO-8601 y en UTC, así que 03:02Z son las 00:02 en Buenos Aires. Están en \texttt{docs/contrato} y un test unitario los compara con lo que genera el simulador, así que no se desactualizan. Los decimales largos de posición y progreso salen así del simulador.

\textbf{Telemetría}

\noindent
\begin{minipage}[t]{0.49\textwidth}
\small Patrullando por el pasillo A
\lstinputlisting[basicstyle=\ttfamily\scriptsize]{../contrato/telemetria-patrullando.json}
\end{minipage}\hfill
\begin{minipage}[t]{0.49\textwidth}
\small Batería por debajo del 25~\%
\lstinputlisting[basicstyle=\ttfamily\scriptsize]{../contrato/telemetria-bateria-baja.json}
\end{minipage}

\textbf{Orden y respuestas}

El primer intento de arranque fue con 20~\% de batería y el robot lo rechazó. Después de cargar, el segundo se aceptó.

\noindent
\begin{minipage}[t]{0.49\textwidth}
\lstinputlisting{../contrato/orden-iniciar-mision.json}
\end{minipage}\hfill
\begin{minipage}[t]{0.49\textwidth}
\lstinputlisting{../contrato/ack-aceptada.json}
\lstinputlisting{../contrato/ack-rechazada-bateria-baja.json}
\end{minipage}

\textbf{Alertas y decisión del guardia}

Ante el obstáculo el guardia eligió esperar. La orden para resolver el incidente lleva el id de la alerta y la decisión.

\noindent
\begin{minipage}[t]{0.49\textwidth}
\lstinputlisting{../contrato/alerta-obstaculo.json}
\lstinputlisting{../contrato/orden-resolver-incidente.json}
\end{minipage}\hfill
\begin{minipage}[t]{0.49\textwidth}
\lstinputlisting{../contrato/alerta-persona-zona-restringida.json}
\end{minipage}

Si pasan 5~segundos sin un \texttt{command\_ack}, la app marca la orden como sin respuesta. No se reintenta sola: el guardia decide si la vuelve a mandar.

\section{Escenarios y riesgos}
\label{sec:escenarios}

\subsection{Los cinco escenarios de la demo}

Todos se pueden reproducir desde el panel de simulación de la app, sin robot ni internet.

\begin{enumerate}
  \item \textbf{Ronda normal.} El guardia inicia la ronda completa. El robot pasa por los ocho puntos, el contador de verificados sube de a uno y al final vuelve a la base. Queda registrada en el historial.
  \item \textbf{Batería baja.} Con la batería en 24~\% aparece una advertencia con el botón ``Volver a la base'' a mano. Debajo del 15~\% el robot pide volver por su cuenta y la app no deja reanudar.
  \item \textbf{Desconexión.} Se corta la conexión desde el panel. Los datos quedan en gris con su antigüedad, la orden que se mande queda sin respuesta y, al reconectar, todo se pone al día. Si el corte supera el minuto, el robot vuelve solo a la base.
  \item \textbf{Incidente.} Una silla bloquea el pasillo B. En la demo el obstáculo se provoca desde el panel con el robot en ese pasillo, entre CP7 y CP8, porque aparece en la zona donde está el robot. El robot frena, pasa a incidente y el guardia elige: esperar a que alguien la saque, saltear ese punto o mandarlo a la base.
  \item \textbf{Persona en zona restringida.} Son las 3 de la mañana y el robot detecta a alguien en el laboratorio. Se detiene a distancia, avisa por el parlante y el guardia decide si es un intruso (y sigue el procedimiento de seguridad), un empleado autorizado o una falsa alarma. La decisión queda en el historial.
\end{enumerate}

Las capturas de cada escenario están en \texttt{docs/evidencias/} del repositorio. El video de demostración muestra un flujo normal y uno con falla: \pendiente{link al video de demostración}

\subsection{Riesgos y medidas}

\begin{table}[H]
\centering\small
\begin{tabularx}{\textwidth}{@{}L L@{}}
\toprule
\textbf{Riesgo} & \textbf{Medida} \\
\midrule
El robot se acerca demasiado a una persona o la asusta & Velocidad tope de 0,8~m/s, se detiene a 2~m de cualquier persona y le da paso. En zonas restringidas avisa por el parlante en lugar de acercarse \\
El guardia cancela o detiene por error & Cancelar pide confirmación. Detener exige mantener apretado. Los dos botones están separados del resto y tienen otro color \\
Se pierde la conexión y el guardia sigue mirando datos viejos & Los datos se marcan como antiguos con su edad, las órdenes muestran ``sin respuesta'' y el robot vuelve solo a la base si el corte supera el minuto \\
Se vulnera la privacidad de quien aparece en las rondas & No hay reconocimiento facial ni se guardan imágenes. El historial guarda zona, hora y la clasificación del guardia, nada más \\
El guardia se cansa de falsas alarmas y deja de prestarles atención & Las alertas tienen severidad y hay que reconocerlas. Las clasificadas como falsa alarma quedan registradas para ajustar la ruta o el sensor \\
\bottomrule
\end{tabularx}
\end{table}

\section{Arquitectura del APK}
\label{sec:arquitectura}

La app está hecha con Kotlin y Jetpack Compose, siguiendo MVVM. El código se divide en cuatro capas: \texttt{ui}, \texttt{presentation}, \texttt{domain} y \texttt{data}. La UI observa el \texttt{StateFlow} que expone cada ViewModel y le manda eventos. Nunca toca el simulador.

El dominio no depende de ninguna otra capa. Tiene el modelo, las reglas y las interfaces de los repositorios. La capa de datos implementa esas interfaces. \texttt{RobotRepository} es la única puerta a la telemetría, las órdenes y los eventos del robot. El día que exista un gateway real, se escribe otra implementación de esa interfaz y el resto de la app no cambia.

\begin{figure}[H]
\centering
% Diagrama de capas. Lo usan el informe (secciones/07-arquitectura.tex) y arquitectura.tex
% Necesita tikz con arrows.meta, fit y backgrounds, y los colores noche, fondo, gris, linea y ambar
% Letra fija de 7 pt para que mida lo mismo en el informe (11 pt) y en la exportación
\begin{tikzpicture}[
  font=\fontsize{7}{8.5}\selectfont,
  caja/.style={draw=noche,fill=white,rounded corners=2pt,thick,text width=#1,align=center,
    minimum height=1.2cm,inner sep=3pt,execute at begin node={\hyphenpenalty=10000\relax}},
  caja/.default=3.3cm,
  capa/.style={draw=linea,fill=fondo,rounded corners=4pt,inner xsep=6pt,inner ysep=5pt},
  nombre/.style={font=\fontsize{9}{11}\selectfont\bfseries\color{noche},anchor=south west,inner sep=1pt},
  flecha/.style={-{Stealth[length=2mm]},thick,gris},
  json/.style={{Stealth[length=2mm]}-{Stealth[length=2mm]},thick,ambar},
  crea/.style={-{Stealth[length=2mm]},thick,gris,dashed}]

  % ui
  \node[caja=4.58cm] (screens) at (0.64,0) {\textbf{screens}\\Supervisión, Misiones, Alertas, Historial, panel de simulación, selección de robot};
  \node[caja=4.58cm] (components) at (5.78,0) {\textbf{components}\\Mapa, tarjeta de telemetría, chips de estado, diálogos, botón de detener};
  \node[caja=4.58cm] (navigation) at (10.91,0) {\textbf{navigation}\\AppNavHost y rutas};

  % presentation
  \node[caja=4.58cm] (vms) at (0.64,-2.2) {\textbf{Un ViewModel y su UiState por pantalla}\\RobotSelect, Supervision, Missions, Alerts, History};
  \node[caja=4.58cm] (mainvm) at (5.78,-2.2) {\textbf{MainViewModel}\\Badge de alertas sin reconocer};
  \node[caja=4.58cm] (simvm) at (10.91,-2.2) {\textbf{SimulationViewModel}\\Panel de simulación};

  % domain
  \node[caja] (model) at (0,-4.4) {\textbf{model}\\Telemetría, estados, misiones, alertas, órdenes};
  \node[caja] (rules) at (3.85,-4.4) {\textbf{rules}\\CommandRules, BatteryRules, Freshness, IncidentRules};
  \node[caja] (repos) at (7.7,-4.4) {\textbf{repository} (interfaces)\\RobotRepository, HistoryRepository, SimulationControls};
  \node[caja] (usecase) at (11.55,-4.4) {\textbf{usecase}\\HistoryRecorder: pasa misiones y eventos al historial};

  % data
  \node[caja] (contract) at (0,-6.8) {\textbf{contract}\\DTOs, mappers y ContractJson};
  \node[caja] (simrepo) at (3.85,-6.8) {\textbf{sim}\\SimulatedRobotRepository};
  \node[caja] (roomrepo) at (7.7,-6.8) {\textbf{local}\\RoomHistoryRepository};
  \node[caja] (db) at (11.55,-6.8) {\textbf{local}\\VigiaDatabase (Room): rondas y eventos};
  \node[caja] (sim) at (3.85,-8.8) {\textbf{RobotSimulator}\\Lado robot, tick de 1~s};

  \begin{scope}[on background layer]
    \node[capa,fit=(screens)(components)(navigation)] (ui) {};
    \node[capa,fit=(vms)(mainvm)(simvm)] (presentation) {};
    \node[capa,fit=(model)(rules)(repos)(usecase)] (domain) {};
    \node[capa,fit=(contract)(simrepo)(roomrepo)(db)(sim)] (data) {};
  \end{scope}
  \node[nombre] at (ui.north west) {ui};
  \node[nombre] at (presentation.north west) {presentation};
  \node[nombre] at (domain.north west) {domain};
  \node[nombre] at (data.north west) {data};

  % di
  \node[caja=8.4cm,minimum height=0.9cm] (di) at (5.78,-11.0) {\textbf{di/AppContainer}\\Crea el simulador, la base y los repositorios, arranca HistoryRecorder y le da las dependencias a las factories de los ViewModels};

  % Dependencias entre capas: cada flecha va de quien usa a quien es usado
  \draw[flecha] ([xshift=-1.5cm]ui.south) -- node[right]{eventos} ([xshift=-1.5cm]presentation.north);
  \draw[flecha] ([xshift=1.5cm]presentation.north) -- node[right]{StateFlow\textless UiState\textgreater} ([xshift=1.5cm]ui.south);
  \draw[flecha] (presentation.south) -- node[right]{modelo, reglas e interfaces} (domain.north);
  \draw[flecha] (simrepo.north) -- node[left,pos=0.4]{implementa} (repos.south west);
  \draw[flecha] (roomrepo.north) -- node[right,pos=0.4]{implementa} (repos.south);
  \draw[flecha] (usecase) -- (repos);
  \draw[flecha] (roomrepo) -- (db);
  \draw[flecha] (simrepo) -- (contract);
  \draw[json] (simrepo) -- node[fill=fondo,inner sep=1pt]{JSON del contrato} (sim);

  \draw[crea] (di.north) -- node[right]{crea} (data.south);
  \draw[crea] (di.east) -| ([xshift=3mm]presentation.east) -- (presentation.east);
\end{tikzpicture}

\caption{Capas de la app. Las flechas llenas indican quién conoce a quién; las punteadas, quién crea cada pieza.}
\label{fig:arquitectura}
\end{figure}

\begin{lstlisting}
com.vigia
 |- domain/model       telemetría, estados, misiones, alertas, órdenes, plano
 |- domain/rules       CommandRules, BatteryRules, Freshness, IncidentRules
 |- domain/repository  RobotRepository, HistoryRepository, SimulationControls
 |- domain/usecase     HistoryRecorder: pasa misiones y eventos al historial
 |- data/contract      DTOs del contrato JSON, mappers y ContractJson
 |- data/sim           RobotSimulator y SimulatedRobotRepository
 |- data/local         Room: VigiaDatabase, entidades, DAOs y RoomHistoryRepository
 |- presentation       un ViewModel y su UiState por pantalla, MainViewModel
 |- ui                 pantallas, componentes, tema y navegación
 |- di                 AppContainer (inyección manual)
\end{lstlisting}

Cada pantalla tiene su ViewModel: selección de robot, supervisión, misiones, alertas, historial y panel de simulación. \texttt{MainViewModel} es el de la navegación y solo cuenta las alertas sin reconocer para mostrar el badge. \texttt{SimulationViewModel} usa \texttt{SimulationControls}, una interfaz aparte para que el resto de las pantallas no pueda tocar la simulación.

\texttt{AppContainer} crea una sola vez el simulador, la base y los repositorios. También arranca \texttt{HistoryRecorder}, que escucha al robot y guarda cada ronda y cada evento en el historial. Las factories de los ViewModels toman sus dependencias de ahí.

\subsection{Simulador}

El simulador tiene dos partes. \texttt{RobotSimulator} hace de robot: recibe órdenes en JSON, contesta con un \texttt{command\_ack} y genera telemetría y alertas en JSON. No usa corrutinas y el tiempo solo avanza cuando se lo piden, así que en los tests es determinístico. \texttt{SimulatedRobotRepository} hace de app: serializa cada orden, parsea lo que vuelve y lo pasa al dominio. Entre los dos solo viaja el JSON del contrato (sección~\ref{sec:contrato}).

El repositorio simulado corre una corrutina con un tick de un segundo y un multiplicador (x1 o x5) para acortar las demos frente al público. En cada tick el simulador mueve el robot sobre la ruta a 0,8~m/s, descuenta batería, se detiene unos segundos en cada punto de control y lo marca como verificado. Desde el panel de simulación se le puede cortar la conexión, pedir que rechace la próxima orden, bajar la batería o provocar un obstáculo, una puerta abierta o una persona en zona restringida. Nada de esto necesita internet.

\subsection{Persistencia}

El historial de rondas y eventos se guarda en Room. Elegimos Room porque son registros con estructura fija que se consultan por fecha y por ronda, y porque ya lo usamos en la materia. El historial no guarda imágenes ni datos de personas (ver la sección~\ref{sec:escenarios}).

\section{Decisiones y limitaciones}
\label{sec:decisiones}

\begin{table}[H]
\centering\small
\begin{tabularx}{\textwidth}{@{}L L@{}}
\toprule
\textbf{Decisión} & \textbf{Por qué} \\
\midrule
Órdenes discretas en lugar de joystick & El guardia supervisa una ronda, no maneja un vehículo. Un joystick invita a manejar a distancia y suma riesgo \\
Tema oscuro por defecto & La central está en penumbra durante toda la jornada \\
Detener con pulsación sostenida & Es la orden que más conviene que llegue y la que menos conviene mandar por un toque accidental \\
El repositorio simulado habla JSON & Así el contrato se prueba en la app y no queda solo en el documento \\
Inyección manual con un \texttt{AppContainer} & Para el tamaño de la app alcanza; Hilt sumaría configuración sin aportar \\
Sin imágenes en el historial & Privacidad de las personas que aparecen en las rondas \\
\bottomrule
\end{tabularx}
\end{table}

\subsection{Limitaciones}

\begin{itemize}
  \item No hay robot ni gateway: todo el comportamiento lo genera el simulador, con reglas simples y deterministas.
  \item La detección de personas y de puertas abiertas se simula. En un robot real haría falta un modelo de visión propio.
  \item El plano es una simplificación de una planta baja; no corresponde a la planta real de la sede.
  \item Hay un solo operador. No contemplamos dos guardias mandando órdenes al mismo robot.
  \item \pendiente{completar con lo que surja durante la implementación}
\end{itemize}

\section{Fuentes}
\label{sec:fuentes}

\begin{itemize}
  \item Unitree Robotics, página del producto Go2: \url{https://www.unitree.com/go2}
  \item Unitree Robotics, documentación para desarrolladores: \url{https://support.unitree.com/home/en/developer}
  \item Consigna del trabajo práctico: \texttt{docs/consigna/TP-HMI-Ottoman.pdf}
  \item Hyundai Motor Group, lanzamiento del Factory Safety Service Robot con Boston Dynamics (2021): \url{https://www.hyundai.com/au/en/news/safety/hyundai-motor-group-launches-Factory-safety-service-robot-first-project-with-boston-dynamics}
  \item Corps Security, demostración de patrullaje con robots en un campus de Londres: \url{https://www.corpssecurity.co.uk/case-study/security-monitoring/robotic-patrols-demo-at-diverse-media-client-site/}
\end{itemize}


\end{document}
